% =====================================================================
% Sección: Pregunta de investigación, hipótesis y diseño experimental
% Responsable: Parte 1 — David Fernando Benjumea Mora
% =====================================================================
\section{Pregunta de investigación e hipótesis}
\label{sec:pregunta}

\subsection{Método científico aplicado}

El trabajo sigue las etapas del método científico \cite{bunge1995ciencia}. Este método obliga a fijar la pregunta, la hipótesis y el criterio de
decisión \emph{antes} de mirar los datos, lo que evita acomodar las conclusiones a los
resultados. La tabla~\ref{tab:metodo} resume cómo se concreta cada etapa.

\begin{table}[h]
\centering
\small
\begin{tabular}{p{2.6cm}p{9.4cm}}
\toprule
\textbf{Etapa} & \textbf{En este trabajo} \\
\midrule
Observación & En una MANET jerárquica la movilidad ocurre en dos niveles (clúster y
nodo), que afectan enlaces de naturaleza distinta (backbone y enlaces locales). \\
Pregunta & ¿Cuál de los dos niveles de movilidad degrada más el desempeño, y cómo
depende de la velocidad? (sección~\ref{sec:preg}) \\
Hipótesis & La movilidad de clúster degrada más que la de nodo
(sección~\ref{sec:hip}). \\
Predicción & Con la misma velocidad, la PDR del escenario con solo movilidad de clúster
será menor y su retardo mayor que en el escenario con solo movilidad de nodo. \\
Experimento & Simulación controlada en ns-3 con cuatro escenarios, cuatro velocidades y
diez réplicas por configuración (sección~\ref{sec:diseno-exp}). \\
Análisis & Medias, intervalos de confianza del 95\,\% y pruebas $t$ de Welch
(sección~\ref{sec:analisis-est}). \\
Conclusión & Se acepta o se rechaza la hipótesis según la regla de decisión
definida de antemano. \\
\bottomrule
\end{tabular}
\caption{Etapas del método científico aplicadas al estudio.}
\label{tab:metodo}
\end{table}

\subsection{Pregunta de investigación}
\label{sec:preg}

\begin{quote}
\emph{En una MANET jerárquica de dos niveles con enrutamiento AODV, ¿qué degrada más el
desempeño de la red —la movilidad a nivel de clúster o la movilidad a nivel de nodo— y
cómo cambia ese efecto al aumentar la velocidad?}
\end{quote}

\subsection{Hipótesis}
\label{sec:hip}

\textbf{Hipótesis principal.} A una misma velocidad, la movilidad de clúster degrada más
el desempeño que la movilidad de nodo: produce una menor tasa de entrega de paquetes y un
mayor retardo. La razón esperada es que el movimiento de un clúster rompe los enlaces
entre cluster heads, por los que pasa todo el tráfico entre clústeres, mientras que el
movimiento de los nodos solo altera rutas cortas dentro de cada clúster.

Sea $\mu_{k}^{\mathrm{PDR}}(v)$ la PDR media del escenario $k$ a la velocidad $v$, y
$\mu_{k}^{D}(v)$ su retardo medio. Para cada velocidad $v$ se plantean:

\begin{align*}
H_0^{\mathrm{PDR}}&: \mu_{E2}^{\mathrm{PDR}}(v) = \mu_{E1}^{\mathrm{PDR}}(v) &
H_1^{\mathrm{PDR}}&: \mu_{E2}^{\mathrm{PDR}}(v) < \mu_{E1}^{\mathrm{PDR}}(v) \\
H_0^{D}&: \mu_{E2}^{D}(v) = \mu_{E1}^{D}(v) &
H_1^{D}&: \mu_{E2}^{D}(v) > \mu_{E1}^{D}(v)
\end{align*}

donde E1 es el escenario con solo movilidad de nodo y E2 el escenario con solo
movilidad de clúster (tabla~\ref{tab:escenarios}).

\textbf{Hipótesis secundarias.} La hipótesis principal compara los dos tipos de
movilidad a una velocidad dada; las secundarias consideran el efecto de la velocidad y de la combinación de ambos movimientos:
\begin{itemize}
    \item \textbf{HS1:} en los escenarios con movimiento (E1, E2 y E3), el desempeño
    empeora a medida que aumenta la velocidad.
    \item \textbf{HS2:} el escenario combinado E3 tiene un desempeño igual o peor que el
    peor de E1 y E2.
    \item \textbf{HS3:} la diferencia entre E1 y E2 crece con la velocidad.
\end{itemize}

\subsection{Variables}

\begin{table}[h]
\centering
\small
\begin{tabular}{p{2.6cm}p{9.4cm}}
\toprule
\textbf{Tipo} & \textbf{Variables} \\
\midrule
Independientes & Tipo de movilidad (ninguna, nodo, clúster, ambas) y velocidad
$v \in \{1, 5, 10, 20\}$\,m/s. \\
Dependientes & PDR, retardo extremo a extremo medio, throughput y jitter medio. \\
Controladas & Número y agrupación de nodos, geometría, estándar WiFi, tasa de
transmisión, potencia, modelo de propagación, protocolo de enrutamiento, tráfico y
tiempo de simulación (tabla~\ref{tab:parametros}). \\
\bottomrule
\end{tabular}
\caption{Variables del experimento.}
\label{tab:variables}
\end{table}

{\sloppy
Las métricas se calculan a partir de los contadores que entrega FlowMonitor
(\texttt{txPackets}, \texttt{rxPackets}, \texttt{rxBytes}, \texttt{delaySum},
\texttt{jitterSum}, \texttt{timeFirstTxPacket}, \texttt{timeLastRxPacket}), sumando sobre
todos los flujos $f$ de una corrida:\par}

\begin{align*}
\mathrm{PDR} &= \frac{\sum_f \mathrm{rxPackets}_f}{\sum_f \mathrm{txPackets}_f},
\qquad
\bar{D} = \frac{\sum_f \mathrm{delaySum}_f}{\sum_f \mathrm{rxPackets}_f}, \\[4pt]
\mathrm{Throughput} &= \frac{8\sum_f \mathrm{rxBytes}_f}{t_{\mathrm{fin}} - t_{\mathrm{inicio}}},
\qquad
\bar{J} = \frac{\sum_f \mathrm{jitterSum}_f}{\sum_f (\mathrm{rxPackets}_f - 1)},
\end{align*}

donde $t_{\mathrm{inicio}}$ y $t_{\mathrm{fin}}$ son el inicio y el fin del periodo de
tráfico.

\subsection{Diseño experimental}
\label{sec:diseno-exp}

\begin{table}[h]
\centering
\small
\begin{tabular}{lccp{6.2cm}}
\toprule
\textbf{Escenario} & \textbf{Clúster} & \textbf{Nodo} & \textbf{Propósito} \\
\midrule
E0 & No & No & Control (línea base sin movilidad). \\
E1 & No & Sí & Efecto aislado de la movilidad de nodo. \\
E2 & Sí & No & Efecto aislado de la movilidad de clúster. \\
E3 & Sí & Sí & Efecto combinado de ambos niveles. \\
\bottomrule
\end{tabular}
\caption{Escenarios de movilidad.}
\label{tab:escenarios}
\end{table}

E0 cumple el papel de control: muestra cómo se comporta la red cuando nada se mueve y
sirve de referencia para medir la degradación de los demás escenarios. Los escenarios
E1, E2 y E3 se evalúan a las cuatro velocidades, mientras que E0 se evalúa una sola vez
porque no tiene movimiento. Cada configuración se repite $n = 10$ veces con semillas
distintas (parámetro \texttt{RngRun} de ns-3), lo que da
$(3 \times 4 + 1) \times 10 = 130$ corridas.

Para que E1 y E2 sean comparables, en ambos se mueve \textbf{una sola capa} a la misma
velocidad $v$: los nodos en E1 y los clústeres en E2. Así, la diferencia entre ellos puede atribuirse al tipo de movilidad. Los superclústeres permanecen fijos
en todos los escenarios; su función es organizar la jerarquía y delimitar la región en la
que se mueven sus clústeres.

El alcance de cada canal se fija explícitamente con un modelo de disco
(\texttt{Range\allowbreak{}Propagation\allowbreak{}LossModel}): todo receptor dentro del alcance recibe la señal y
ninguno fuera de él. Con el modelo por defecto de ns-3 (pérdida log-distancia), la
potencia de transmisión por defecto y el umbral de detección de preámbulo de ns-3.43
($-82$\,dBm), un enlace 802.11b a 1\,Mbps solo alcanza unos 50\,m (medido: 10 de 10
paquetes a 50\,m y ninguno a 55\,m), lo que impediría tener un backbone entre clústeres
separados 150\,m. El modelo de disco, además, hace que la geometría del experimento sea
fácil de controlar y de explicar. Un alcance de 60\,m dentro del clúster (comparable
a la región de $\pm 40$\,m en que se mueven los nodos) y de 250\,m entre CH (comparable a
la separación de 150\,m entre clústeres y a su región de $\pm 150$\,m) hace que ambos
tipos de movilidad puedan romper enlaces, que es justamente lo que se quiere medir.

\begin{table}[h]
\centering
\small
\begin{tabular}{p{4.2cm}p{7.8cm}}
\toprule
\textbf{Parámetro} & \textbf{Valor} \\
\midrule
Simulador & ns-3.43 (entorno de referencia: Ubuntu 22.04; también ejecutable en Windows nativo u otras distribuciones de Linux) \\
Nodos & 20: 4 clústeres de 5 nodos (1 CH fijo + 4 miembros) \\
Jerarquía & Nivel 1: clústeres A, B, C, D; nivel 2: X = \{A, B\}, Y = \{C, D\} \\
Canales & Un canal WiFi por clúster y un canal backbone compartido por los 4 CH \\
Estándar WiFi & IEEE 802.11b, modo ad hoc, tasa fija DSSS 1\,Mbps
(\texttt{ConstantRateWifiManager}) \\
Propagación & Modelo de disco (\texttt{Range\allowbreak{}Propagation\allowbreak{}LossModel}) con retardo de
velocidad constante \\
Alcance efectivo & 60\,m dentro de cada clúster y 250\,m en el backbone entre CH \\
Potencia de transmisión & Valor por defecto de ns-3 \\
Enrutamiento & AODV con parámetros por defecto \\
Tráfico & 8 flujos UDP \texttt{OnOff} de tasa constante, paquetes de 512\,B a
4\,paquetes/s (16\,kbit/s por flujo) \\
Pares de tráfico & Dentro de cada superclúster: A$\rightarrow$B, B$\rightarrow$A,
C$\rightarrow$D, D$\rightarrow$C; entre superclústeres: A$\rightarrow$C,
B$\rightarrow$D, C$\rightarrow$A, D$\rightarrow$B (miembro a miembro) \\
Tiempos & Simulación de 200\,s; tráfico desde un instante uniforme en $[10, 11]$\,s
hasta 195\,s \\
Movilidad de clúster & Random Waypoint relativo al centro de su superclúster, región de
$\pm 150$\,m, velocidad constante $v$, pausa de 2\,s \\
Movilidad de nodo & Random Walk 2D relativo al centro de su clúster, región de
$\pm 40$\,m, velocidad constante $v$ \\
Posiciones iniciales & Superclústeres en $(350, 300)$ y $(650, 300)$\,m; clústeres en
$x = 275, 425, 575, 725$\,m sobre $y = 300$\,m; CH en el centro y miembros en un anillo de
25\,m de radio \\
\bottomrule
\end{tabular}
\caption{Parámetros constantes de la simulación. Los valores geométricos son valores de
partida: si en E0 la PDR resulta menor al 95\,\%, se ajustan para garantizar una red
conectada y el cambio se documenta.}
\label{tab:parametros}
\end{table}

\subsection{Método de análisis estadístico}
\label{sec:analisis-est}

Cada corrida es una realización independiente del proceso estocástico de movilidad. Para cada configuración y cada métrica se calcula la
media muestral $\bar{x}$, la desviación estándar $s$ y el intervalo de confianza del
95\,\%:
\[
\bar{x} \pm t_{0.975,\,n-1}\,\frac{s}{\sqrt{n}},
\qquad t_{0.975,\,9} \approx 2.262 .
\]

La hipótesis principal se contrasta, para cada velocidad, con una prueba $t$ de Welch
unilateral entre E2 y E1 con nivel de significancia $\alpha = 0.05$. Se define de
antemano la siguiente \textbf{regla de decisión}: la hipótesis principal se considera
respaldada si se rechaza $H_0^{\mathrm{PDR}}$ en al menos tres de las cuatro
velocidades; si se rechaza en menos, no se acepta. El mismo criterio se aplica al
retardo. Las hipótesis secundarias se evalúan comparando las medias y sus intervalos de
confianza en las gráficas de desempeño contra velocidad.
